% Options for packages loaded elsewhere
\PassOptionsToPackage{unicode}{hyperref}
\PassOptionsToPackage{hyphens}{url}
\documentclass[
  11pt,
  oneside]{report}
\usepackage{xcolor}
\usepackage[margin=1in]{geometry}
\usepackage{amsmath,amssymb}
\setcounter{secnumdepth}{-\maxdimen} % remove section numbering
\usepackage{iftex}
\ifPDFTeX
  \usepackage[T1]{fontenc}
  \usepackage[utf8]{inputenc}
  \usepackage{textcomp} % provide euro and other symbols
\else % if luatex or xetex
  \usepackage{unicode-math} % this also loads fontspec
  \defaultfontfeatures{Scale=MatchLowercase}
  \defaultfontfeatures[\rmfamily]{Ligatures=TeX,Scale=1}
\fi
\usepackage{lmodern}
\ifPDFTeX\else
  % xetex/luatex font selection
\fi
% Use upquote if available, for straight quotes in verbatim environments
\IfFileExists{upquote.sty}{\usepackage{upquote}}{}
\IfFileExists{microtype.sty}{% use microtype if available
  \usepackage[]{microtype}
  \UseMicrotypeSet[protrusion]{basicmath} % disable protrusion for tt fonts
}{}
\makeatletter
\@ifundefined{KOMAClassName}{% if non-KOMA class
  \IfFileExists{parskip.sty}{%
    \usepackage{parskip}
  }{% else
    \setlength{\parindent}{0pt}
    \setlength{\parskip}{6pt plus 2pt minus 1pt}}
}{% if KOMA class
  \KOMAoptions{parskip=half}}
\makeatother
\usepackage{color}
\usepackage{fancyvrb}
\newcommand{\VerbBar}{|}
\newcommand{\VERB}{\Verb[commandchars=\\\{\}]}
\DefineVerbatimEnvironment{Highlighting}{Verbatim}{commandchars=\\\{\}}
% Add ',fontsize=\small' for more characters per line
\usepackage{framed}
\definecolor{shadecolor}{RGB}{248,248,248}
\newenvironment{Shaded}{\begin{snugshade}}{\end{snugshade}}
\newcommand{\AlertTok}[1]{\textcolor[rgb]{0.94,0.16,0.16}{#1}}
\newcommand{\AnnotationTok}[1]{\textcolor[rgb]{0.56,0.35,0.01}{\textbf{\textit{#1}}}}
\newcommand{\AttributeTok}[1]{\textcolor[rgb]{0.13,0.29,0.53}{#1}}
\newcommand{\BaseNTok}[1]{\textcolor[rgb]{0.00,0.00,0.81}{#1}}
\newcommand{\BuiltInTok}[1]{#1}
\newcommand{\CharTok}[1]{\textcolor[rgb]{0.31,0.60,0.02}{#1}}
\newcommand{\CommentTok}[1]{\textcolor[rgb]{0.56,0.35,0.01}{\textit{#1}}}
\newcommand{\CommentVarTok}[1]{\textcolor[rgb]{0.56,0.35,0.01}{\textbf{\textit{#1}}}}
\newcommand{\ConstantTok}[1]{\textcolor[rgb]{0.56,0.35,0.01}{#1}}
\newcommand{\ControlFlowTok}[1]{\textcolor[rgb]{0.13,0.29,0.53}{\textbf{#1}}}
\newcommand{\DataTypeTok}[1]{\textcolor[rgb]{0.13,0.29,0.53}{#1}}
\newcommand{\DecValTok}[1]{\textcolor[rgb]{0.00,0.00,0.81}{#1}}
\newcommand{\DocumentationTok}[1]{\textcolor[rgb]{0.56,0.35,0.01}{\textbf{\textit{#1}}}}
\newcommand{\ErrorTok}[1]{\textcolor[rgb]{0.64,0.00,0.00}{\textbf{#1}}}
\newcommand{\ExtensionTok}[1]{#1}
\newcommand{\FloatTok}[1]{\textcolor[rgb]{0.00,0.00,0.81}{#1}}
\newcommand{\FunctionTok}[1]{\textcolor[rgb]{0.13,0.29,0.53}{\textbf{#1}}}
\newcommand{\ImportTok}[1]{#1}
\newcommand{\InformationTok}[1]{\textcolor[rgb]{0.56,0.35,0.01}{\textbf{\textit{#1}}}}
\newcommand{\KeywordTok}[1]{\textcolor[rgb]{0.13,0.29,0.53}{\textbf{#1}}}
\newcommand{\NormalTok}[1]{#1}
\newcommand{\OperatorTok}[1]{\textcolor[rgb]{0.81,0.36,0.00}{\textbf{#1}}}
\newcommand{\OtherTok}[1]{\textcolor[rgb]{0.56,0.35,0.01}{#1}}
\newcommand{\PreprocessorTok}[1]{\textcolor[rgb]{0.56,0.35,0.01}{\textit{#1}}}
\newcommand{\RegionMarkerTok}[1]{#1}
\newcommand{\SpecialCharTok}[1]{\textcolor[rgb]{0.81,0.36,0.00}{\textbf{#1}}}
\newcommand{\SpecialStringTok}[1]{\textcolor[rgb]{0.31,0.60,0.02}{#1}}
\newcommand{\StringTok}[1]{\textcolor[rgb]{0.31,0.60,0.02}{#1}}
\newcommand{\VariableTok}[1]{\textcolor[rgb]{0.00,0.00,0.00}{#1}}
\newcommand{\VerbatimStringTok}[1]{\textcolor[rgb]{0.31,0.60,0.02}{#1}}
\newcommand{\WarningTok}[1]{\textcolor[rgb]{0.56,0.35,0.01}{\textbf{\textit{#1}}}}
\usepackage{longtable,booktabs,array}
\usepackage{calc} % for calculating minipage widths
% Correct order of tables after \paragraph or \subparagraph
\usepackage{etoolbox}
\makeatletter
\patchcmd\longtable{\par}{\if@noskipsec\mbox{}\fi\par}{}{}
\makeatother
% Allow footnotes in longtable head/foot
\IfFileExists{footnotehyper.sty}{\usepackage{footnotehyper}}{\usepackage{footnote}}
\makesavenoteenv{longtable}
\setlength{\emergencystretch}{3em} % prevent overfull lines
\providecommand{\tightlist}{%
  \setlength{\itemsep}{0pt}\setlength{\parskip}{0pt}}
\usepackage{bookmark}
\IfFileExists{xurl.sty}{\usepackage{xurl}}{} % add URL line breaks if available
\urlstyle{same}
\hypersetup{
  hidelinks,
  pdfcreator={LaTeX via pandoc}}

\author{}
\date{}

\begin{document}

\chapter{The Renoise Theme-to-Screenshot Preview
System}\label{the-renoise-theme-to-screenshot-preview-system}

\section{A Pixel-Accurate Theme Preview Engine for Closed-Source
Software}\label{a-pixel-accurate-theme-preview-engine-for-closed-source-software}

\begin{center}\rule{0.5\linewidth}{0.5pt}\end{center}

\textbf{Abstract.} This paper documents a pixel-accurate theme preview
system that renders Renoise Digital Audio Workstation (DAW) color themes
without requiring a Renoise installation. The system maps every pixel in
reference screenshots to semantic UI elements via binary pixel maps,
then recolors those pixels using colors extracted from \texttt{.xrnc}
theme files. The result is a three-pass rendering pipeline that produces
preview images indistinguishable from native screenshots --- at a
fraction of the computational and licensing cost. We describe the
original ``rainbow theme trick,'' the refined ``white-baseline + green
probe'' method, the three-pass rendering algorithm, and the headless
automation pipeline that makes map generation possible at scale.

\begin{center}\rule{0.5\linewidth}{0.5pt}\end{center}

\section{1. Executive Summary}\label{executive-summary}

Renoise is a professional music tracker and Digital Audio Workstation
used by electronic musicians worldwide. Like many DAWs, it supports
custom color themes via \texttt{.xrnc} XML files containing 70+ color
definitions. However, Renoise provides no built-in theme preview
mechanism. Music producers who want to evaluate a theme must install it,
restart Renoise, and manually inspect every view. Theme sharing
communities have historically relied on manual screenshots ---
inconsistent, time-consuming, and impossible to generate at scale.

The Renoise Theme-to-Screenshot Preview System solves this by treating
theme preview as a \textbf{pixel recoloring problem} rather than a UI
rendering problem. Instead of reimplementing Renoise's closed-source
interface (impossible), we:

\begin{enumerate}
\def\labelenumi{\arabic{enumi}.}
\tightlist
\item
  \textbf{Map} every pixel in a reference screenshot to a semantic UI
  element (\texttt{Main\_Back}, \texttt{Body\_Font}, etc.)
\item
  \textbf{Parse} any \texttt{.xrnc} file to extract its color
  definitions
\item
  \textbf{Recolor} the mapped pixels using the theme's colors
\item
  \textbf{Composite} text and anti-aliased details using a special
  black/white text-mask theme
\end{enumerate}

The result is a web application where users upload a \texttt{.xrnc} file
and receive accurate PNG previews of Pattern Editor, Mixer, and Waveform
views within seconds --- no Renoise installation required.

\textbf{Key innovations:} - Binary pixel maps derived from automated
green-probe differential analysis - A three-pass rendering algorithm
that handles anti-aliasing, boundary smoothing, and text compositing - A
headless Xvfb automation pipeline that generates 71 theme variants in
\textasciitilde3 minutes - A live Theme Creator tool with 70 color
pickers ranked by pixel coverage

\begin{center}\rule{0.5\linewidth}{0.5pt}\end{center}

\section{2. The Problem Space}\label{the-problem-space}

\subsection{2.1 What is a Renoise Theme?}\label{what-is-a-renoise-theme}

A Renoise theme is an XML file with the extension \texttt{.xrnc}
(Renoise Color). It defines colors for 70+ UI elements, from the main
application background (\texttt{Main\_Back}) to the color of automation
curve edges (\texttt{Automation\_Line\_Edge}). A minimal theme looks
like this:

\begin{Shaded}
\begin{Highlighting}[]
\FunctionTok{\textless{}?xml}\OtherTok{ version=}\StringTok{"1.0"}\OtherTok{ encoding=}\StringTok{"UTF{-}8"}\FunctionTok{?\textgreater{}}
\NormalTok{\textless{}}\KeywordTok{SkinColors}\OtherTok{ doc\_version=}\StringTok{"12"}\NormalTok{\textgreater{}}
\NormalTok{  \textless{}}\KeywordTok{Main\_Back}\NormalTok{\textgreater{}30,30,30\textless{}/}\KeywordTok{Main\_Back}\NormalTok{\textgreater{}}
\NormalTok{  \textless{}}\KeywordTok{Main\_Font}\NormalTok{\textgreater{}220,220,220\textless{}/}\KeywordTok{Main\_Font}\NormalTok{\textgreater{}}
\NormalTok{  \textless{}}\KeywordTok{Body\_Back}\NormalTok{\textgreater{}40,40,50\textless{}/}\KeywordTok{Body\_Back}\NormalTok{\textgreater{}}
  \CommentTok{\textless{}!{-}{-} ... 70 more elements ... {-}{-}\textgreater{}}
\NormalTok{\textless{}/}\KeywordTok{SkinColors}\NormalTok{\textgreater{}}
\end{Highlighting}
\end{Shaded}

These colors control every visible surface of Renoise's interface: the
pattern grid, mixer channels, device chains, automation envelopes, VU
meters, buttons, sliders, scrollbars, and tooltips. Changing
\texttt{Main\_Back} affects roughly 46\% of the screen. Changing
\texttt{VuMeter\_Peak} affects perhaps 0.01\%.

\subsection{2.2 The Preview Gap}\label{the-preview-gap}

Before this system existed, theme authors and curators faced three
problems:

\begin{enumerate}
\def\labelenumi{\arabic{enumi}.}
\tightlist
\item
  \textbf{No built-in preview.} Renoise does not generate thumbnails or
  previews. To see a theme, you must load it.
\item
  \textbf{Manual screenshots are inconsistent.} Different users run
  different screen resolutions, window sizes, and song projects. One
  user's screenshot might show the pattern editor; another's might show
  the mixer. There is no standardization.
\item
  \textbf{Renoise is required.} Generating a screenshot means having
  Renoise installed, licensed, and running. This creates friction for
  web-based theme galleries and mobile users.
\end{enumerate}

The gap between ``upload a theme'' and ``see how it looks'' was the
single biggest friction point in the Renoise theme-sharing ecosystem.

\subsection{2.3 Why Not Just Render the
UI?}\label{why-not-just-render-the-ui}

The most obvious solution --- parse the theme and draw the UI from
scratch --- is infeasible. Renoise is closed-source proprietary
software. Its UI contains thousands of custom widgets, proprietary bevel
and shading algorithms, textured panels, and anti-aliased text rendered
with unknown font metrics. Reverse-engineering the entire renderer would
be a multi-year project.

We needed a different approach.

\begin{center}\rule{0.5\linewidth}{0.5pt}\end{center}

\section{3. System Architecture
Overview}\label{system-architecture-overview}

The system is a Node.js/Express web application with a clear separation
between upload processing, preview generation, and content delivery.

\begin{verbatim}
┌─────────────────────────────────────────────────────────────────────┐
│                         CLIENT (Browser)                            │
│  ┌─────────────┐  ┌─────────────┐  ┌─────────────────────────────┐  │
│  │ Theme Gallery│  │ Upload Form │  │ Theme Creator (/create)     │  │
│  └──────┬──────┘  └──────┬──────┘  └─────────────┬───────────────┘  │
└─────────┼────────────────┼───────────────────────┼──────────────────┘
          │                │ .xrnc upload          │ color map JSON
          ▼                ▼                       ▼
┌─────────────────────────────────────────────────────────────────────┐
│                      EXPRESS SERVER (Node.js)                       │
│  ┌────────────┐  ┌─────────────┐  ┌─────────────┐  ┌─────────────┐ │
│  │ Multer     │→ │ lib/parser  │→ │lib/categorize│→ │lib/palette  │ │
│  │ (upload)   │  │ (XML parse) │  │ (auto-tags)  │  │ (SVG strip) │ │
│  └────────────┘  └──────┬──────┘  └─────────────┘  └─────────────┘ │
│                         │ elementColorMap                          │
│                         ▼                                            │
│              ┌─────────────────────┐                                 │
│              │ lib/preview-renderer│                                 │
│              │  (3-pass pixel map) │                                 │
│              └──────────┬──────────┘                                 │
│                         │ PNG previews                               │
│                         ▼                                            │
│              ┌─────────────────────┐                                 │
│              │   lib/database      │                                 │
│              │   (SQLite WAL)      │                                 │
│              └─────────────────────┘                                 │
└─────────────────────────────────────────────────────────────────────┘
\end{verbatim}

\subsection{3.1 Upload Processing
Pipeline}\label{upload-processing-pipeline}

When a user uploads a \texttt{.xrnc} file, it flows through five stages:

\begin{enumerate}
\def\labelenumi{\arabic{enumi}.}
\tightlist
\item
  \textbf{Multer} handles file upload, sanitizes the filename with
  \texttt{path.basename()}, and stores it in
  \texttt{public/uploads/themes/}.
\item
  \textbf{\texttt{lib/parser.js}} parses the XML, extracts all color
  values, matches element names against \texttt{ROLE\_RULES} regex
  patterns, and assigns semantic roles (\texttt{background},
  \texttt{text}, \texttt{ui}, \texttt{accent}) with visual weights.
  Crucially, it exports \texttt{elementColorMap} --- a plain object
  mapping each element name to its \texttt{{[}r,\ g,\ b{]}} tuple. This
  is the \textbf{single source of truth} used by every downstream step.
\item
  \textbf{\texttt{lib/categorize.js}} converts colors to HSL and runs
  weighted statistical analysis to auto-tag themes: dark/light/medium,
  neon/pastel/monochrome, warm/cool/mixed, and dominant color families.
\item
  \textbf{\texttt{lib/palette.js}} generates an SVG palette strip by
  tiering colors into MAIN, SECONDARY, UI, and ACCENTS bands based on
  accumulated weight.
\item
  \textbf{\texttt{lib/preview-renderer.js}} renders up to three PNG
  previews (Pattern Editor, Mixer, Waveform) using the binary pixel map
  system described in Section 4.
\end{enumerate}

All of this completes in under 2 seconds for a typical theme.

\subsection{3.2 Database \& Frontend}\label{database-frontend}

Theme metadata, tags, preview paths, and auto-generated statistics are
stored in SQLite with WAL mode enabled for concurrent read performance.
The frontend uses EJS templates, PicoCSS for styling, and vanilla
JavaScript. Preview images on the detail page use a lazy-loading
\texttt{data-src} pattern --- images are only fetched when the user
clicks a preview tab, saving \textasciitilde66\% of bandwidth on initial
page load.

\begin{center}\rule{0.5\linewidth}{0.5pt}\end{center}

\section{4. The Pixel Map System (Core
Innovation)}\label{the-pixel-map-system-core-innovation}

The pixel map system is the heart of this project. It is what makes
pixel-accurate preview possible without access to Renoise's source code.

\subsection{4.1 Why Pixel Maps?}\label{why-pixel-maps}

Think of a pixel map like a paint-by-numbers template. In a
paint-by-numbers kit, every region of the canvas has a number, and a
separate key tells you which color to use for each number. Our system
works the same way:

\begin{itemize}
\tightlist
\item
  The \textbf{pixel map} (\texttt{maps/pattern.bin}) is the numbered
  canvas. Each pixel contains a number from 0 to 73 (the element index)
  or 255 (unmatched).
\item
  The \textbf{theme file} (\texttt{.xrnc}) is the color key. It tells us
  that element \#0 (\texttt{Main\_Back}) should be dark gray, element
  \#1 (\texttt{Main\_Font}) should be light gray, and so on.
\item
  The \textbf{renderer} ``paints'' the canvas by replacing each numbered
  pixel with the corresponding color from the key.
\end{itemize}

This approach has a profound advantage: we never need to understand
\emph{how} Renoise draws its UI. We only need to know \emph{which} UI
element owns \emph{which} pixel. The heavy lifting of anti-aliasing,
bevel shading, and font rendering was already done by Renoise when we
captured the reference screenshots. We simply recolor the result.

\subsection{4.2 The Rainbow Theme Trick (Original
Method)}\label{the-rainbow-theme-trick-original-method}

Our first attempt at building pixel maps used what we call the
\textbf{Rainbow Theme Trick}. The idea is elegant:

\begin{enumerate}
\def\labelenumi{\arabic{enumi}.}
\tightlist
\item
  Generate 4 variant themes. In each variant, every Renoise UI element
  is assigned a unique, wildly distinct color from a different quadrant
  of the color wheel.
\item
  Screenshot the same Renoise view (Pattern Editor, same project, same
  window size) for all 4 variants.
\item
  For each pixel location, look at the 4 screenshots. The element that
  ``wins'' the majority vote at that pixel is the owner.
\end{enumerate}

\begin{verbatim}
Variant A:  Main_Back=Red,   Body_Back=Green,  Button_Back=Blue...
Variant B:  Main_Back=Green, Body_Back=Blue,   Button_Back=Red...
Variant C:  Main_Back=Blue,  Body_Back=Red,    Button_Back=Green...
Variant D:  Main_Back=Yellow,Body_Back=Magenta,Button_Back=Cyan...

Pixel (x,y):  A=Red, B=Red, C=Green, D=Red  →  Majority=Red  →  Main_Back
\end{verbatim}

This method works, but it has drawbacks: - \textbf{Color collision
risk.} If two elements happen to share similar colors across variants,
the vote becomes ambiguous. - \textbf{Complexity.} Four variants mean
four Renoise launches, four screenshots, and a voting algorithm that
must handle ties. - \textbf{Anti-aliasing bleed.} Edge pixels may show
blended colors that match none of the assigned colors.

The rainbow method proved the concept, but we needed something cleaner.

\subsection{4.3 The White-Baseline + Green Probe Method (Refined
Method)}\label{the-white-baseline-green-probe-method-refined-method}

The refined method is simpler, more robust, and produces cleaner maps.
We call it the \textbf{White-Baseline + Green Probe} method.

\subsubsection{Step 1: Generate 71 Variant
Themes}\label{step-1-generate-71-variant-themes}

We generate 71 \texttt{.xrnc} files: - \textbf{1 baseline:} Every
element is set to pure white (\texttt{\#FFFFFF}). - \textbf{70 probes:}
Every element is white \emph{except one}, which is set to bright green
(\texttt{\#00FF00}).

\begin{Shaded}
\begin{Highlighting}[]
\CommentTok{// From tools/generate{-}white{-}variants.js}
\KeywordTok{const}\NormalTok{ WHITE }\OperatorTok{=} \StringTok{\textquotesingle{}255,255,255\textquotesingle{}}\OperatorTok{;}
\KeywordTok{const}\NormalTok{ PROBE }\OperatorTok{=} \StringTok{\textquotesingle{}0,255,0\textquotesingle{}}\OperatorTok{;} \CommentTok{// bright green}

\ControlFlowTok{for}\NormalTok{ (}\KeywordTok{const}\NormalTok{ element }\KeywordTok{of}\NormalTok{ ALL\_ELEMENTS) \{}
  \CommentTok{// This element = green, all others = white}
  \KeywordTok{const}\NormalTok{ color }\OperatorTok{=}\NormalTok{ (el }\OperatorTok{===}\NormalTok{ element) }\OperatorTok{?}\NormalTok{ PROBE }\OperatorTok{:}\NormalTok{ WHITE}\OperatorTok{;}
\NormalTok{\}}
\end{Highlighting}
\end{Shaded}

\subsubsection{Step 2: Capture Screenshots via Headless
Automation}\label{step-2-capture-screenshots-via-headless-automation}

Each variant is loaded into Renoise running inside an Xvfb virtual
display (see Section 7). A demo song (``Soon Soon'' by Hunz) is
preloaded so that all UI elements --- pattern data, devices, automation
envelopes, VU meters --- are visible. We screenshot the Pattern Editor
at 1920×1080.

\subsubsection{Step 3: Diff Against
Baseline}\label{step-3-diff-against-baseline}

For each probe screenshot, we compare every pixel against the white
baseline. If a pixel turned significantly greener, it belongs to that
element.

\begin{Shaded}
\begin{Highlighting}[]
\CommentTok{// From tools/build{-}diff{-}maps.js}
\KeywordTok{function} \FunctionTok{isProbe}\NormalTok{(vr}\OperatorTok{,}\NormalTok{ vg}\OperatorTok{,}\NormalTok{ vb}\OperatorTok{,}\NormalTok{ br}\OperatorTok{,}\NormalTok{ bg}\OperatorTok{,}\NormalTok{ bb) \{}
  \KeywordTok{const}\NormalTok{ dg }\OperatorTok{=}\NormalTok{ vg }\OperatorTok{{-}}\NormalTok{ bg}\OperatorTok{;}   \CommentTok{// green delta}
  \KeywordTok{const}\NormalTok{ dr }\OperatorTok{=}\NormalTok{ vr }\OperatorTok{{-}}\NormalTok{ br}\OperatorTok{;}   \CommentTok{// red delta}
  \KeywordTok{const}\NormalTok{ db }\OperatorTok{=}\NormalTok{ vb }\OperatorTok{{-}}\NormalTok{ bb}\OperatorTok{;}   \CommentTok{// blue delta}
  \CommentTok{// The green channel must jump by \textgreater{}120, AND dominate red/blue by \textgreater{}80}
  \ControlFlowTok{return}\NormalTok{ dg }\OperatorTok{\textgreater{}} \DecValTok{120} \OperatorTok{\&\&}\NormalTok{ dg }\OperatorTok{\textgreater{}}\NormalTok{ dr }\OperatorTok{+} \DecValTok{80} \OperatorTok{\&\&}\NormalTok{ dg }\OperatorTok{\textgreater{}}\NormalTok{ db }\OperatorTok{+} \DecValTok{80}\OperatorTok{;}
\NormalTok{\}}
\end{Highlighting}
\end{Shaded}

This threshold is carefully chosen. It is high enough to ignore JPEG
compression artifacts, monitor gamma variations, and subtle
anti-aliasing fringes. It is low enough to catch real UI pixels even
when Renoise applies slight shading or bevel effects.

\subsubsection{Step 4: Overlap
Resolution}\label{step-4-overlap-resolution}

Some pixels naturally belong to multiple elements. A button sits on top
of a background. Text sits on top of a button. When probes overlap, we
apply a simple but effective rule: \textbf{smaller coverage wins.}

\begin{Shaded}
\begin{Highlighting}[]
\CommentTok{// From tools/build{-}diff{-}maps.js — merge phase}
\KeywordTok{const}\NormalTok{ order }\OperatorTok{=}\NormalTok{ elements}
  \OperatorTok{.}\FunctionTok{map}\NormalTok{((name}\OperatorTok{,}\NormalTok{ i) }\KeywordTok{=\textgreater{}}\NormalTok{ (\{ name}\OperatorTok{,} \DataTypeTok{size}\OperatorTok{:}\NormalTok{ elementMaps[name]}\OperatorTok{.}\AttributeTok{size}\OperatorTok{,} \DataTypeTok{idx}\OperatorTok{:}\NormalTok{ i \}))}
  \OperatorTok{.}\FunctionTok{sort}\NormalTok{((a}\OperatorTok{,}\NormalTok{ b) }\KeywordTok{=\textgreater{}}\NormalTok{ a}\OperatorTok{.}\AttributeTok{size} \OperatorTok{{-}}\NormalTok{ b}\OperatorTok{.}\AttributeTok{size}\NormalTok{)}\OperatorTok{;} \CommentTok{// smallest first}

\ControlFlowTok{for}\NormalTok{ (}\KeywordTok{const}\NormalTok{ \{ name}\OperatorTok{,}\NormalTok{ idx \} }\KeywordTok{of}\NormalTok{ order) \{}
  \KeywordTok{const}\NormalTok{ pixels }\OperatorTok{=}\NormalTok{ elementMaps[name]}\OperatorTok{;}
  \ControlFlowTok{for}\NormalTok{ (}\KeywordTok{const}\NormalTok{ p }\KeywordTok{of}\NormalTok{ pixels) \{}
    \ControlFlowTok{if}\NormalTok{ (pixelMap[p] }\OperatorTok{!==} \DecValTok{255}\NormalTok{) \{}
      \CommentTok{// Overlap! Keep existing (smaller/first) element}
      \ControlFlowTok{continue}\OperatorTok{;}
\NormalTok{    \}}
\NormalTok{    pixelMap[p] }\OperatorTok{=}\NormalTok{ idx}\OperatorTok{;}
\NormalTok{  \}}
\NormalTok{\}}
\end{Highlighting}
\end{Shaded}

Because \texttt{Main\_Back} covers 46\% of the screen and
\texttt{Main\_Font} covers only 2.5\%, the font probe is processed
first. Any pixel claimed by both \texttt{Main\_Font} and
\texttt{Main\_Back} is assigned to \texttt{Main\_Font}. This correctly
models the visual stacking order: text sits on top of backgrounds,
buttons sit on top of panels.

\subsubsection{Advantages Over Rainbow
Voting}\label{advantages-over-rainbow-voting}

\begin{longtable}[]{@{}
  >{\raggedright\arraybackslash}p{(\linewidth - 4\tabcolsep) * \real{0.1509}}
  >{\raggedright\arraybackslash}p{(\linewidth - 4\tabcolsep) * \real{0.2830}}
  >{\raggedright\arraybackslash}p{(\linewidth - 4\tabcolsep) * \real{0.5660}}@{}}
\toprule\noalign{}
\begin{minipage}[b]{\linewidth}\raggedright
Aspect
\end{minipage} & \begin{minipage}[b]{\linewidth}\raggedright
Rainbow Voting
\end{minipage} & \begin{minipage}[b]{\linewidth}\raggedright
White-Baseline + Green Probe
\end{minipage} \\
\midrule\noalign{}
\endhead
\bottomrule\noalign{}
\endlastfoot
Variants needed & 4+ & 71 (1 baseline + 70 probes) \\
Color collision risk & High (voting ambiguity) & Zero (single active
color) \\
Anti-aliasing handling & Complex (blended colors) & Simple (green delta
threshold) \\
Overlap resolution & Voting majority & Explicit smallest-first \\
Automation complexity & Moderate & High (but fully scripted) \\
\end{longtable}

The white-baseline method requires more screenshots, but each screenshot
is trivial to analyze. The result is a pixel map with zero ambiguous
pixels and clean boundary definitions.

\subsection{4.4 The Three-Pass Rendering
Algorithm}\label{the-three-pass-rendering-algorithm}

Once we have a pixel map and a parsed theme, we render the preview in
three passes. This algorithm lives in \texttt{lib/preview-renderer.js}.

\subsubsection{Pass 1: Paint Known
Pixels}\label{pass-1-paint-known-pixels}

We iterate over every pixel in the map. If the map index is not
\texttt{255} (UNMATCHED), we look up the corresponding element color and
write it to the output buffer.

\begin{Shaded}
\begin{Highlighting}[]
\CommentTok{// Pre{-}build O(1) lookup array to avoid string matching in the hot loop}
\KeywordTok{const}\NormalTok{ colorByIndex }\OperatorTok{=} \KeywordTok{new} \BuiltInTok{Array}\NormalTok{(elements}\OperatorTok{.}\AttributeTok{length}\NormalTok{)}\OperatorTok{;}
\ControlFlowTok{for}\NormalTok{ (}\KeywordTok{let}\NormalTok{ ei }\OperatorTok{=} \DecValTok{0}\OperatorTok{;}\NormalTok{ ei }\OperatorTok{\textless{}}\NormalTok{ elements}\OperatorTok{.}\AttributeTok{length}\OperatorTok{;}\NormalTok{ ei}\OperatorTok{++}\NormalTok{) \{}
\NormalTok{  colorByIndex[ei] }\OperatorTok{=}\NormalTok{ colors[elements[ei]] }\OperatorTok{||} \KeywordTok{null}\OperatorTok{;}
\NormalTok{\}}

\CommentTok{// Pass 1: paint pixels assigned to a theme element}
\ControlFlowTok{for}\NormalTok{ (}\KeywordTok{let}\NormalTok{ i }\OperatorTok{=} \DecValTok{0}\OperatorTok{;}\NormalTok{ i }\OperatorTok{\textless{}}\NormalTok{ W }\OperatorTok{*}\NormalTok{ H}\OperatorTok{;}\NormalTok{ i}\OperatorTok{++}\NormalTok{) \{}
  \KeywordTok{const}\NormalTok{ idx }\OperatorTok{=}\NormalTok{ map[i]}\OperatorTok{;}
  \ControlFlowTok{if}\NormalTok{ (idx }\OperatorTok{===}\NormalTok{ UNMATCHED) }\ControlFlowTok{continue}\OperatorTok{;}
  \KeywordTok{const}\NormalTok{ c }\OperatorTok{=}\NormalTok{ colorByIndex[idx]}\OperatorTok{;}
  \ControlFlowTok{if}\NormalTok{ (}\OperatorTok{!}\NormalTok{c) }\ControlFlowTok{continue}\OperatorTok{;}
  \KeywordTok{const}\NormalTok{ px }\OperatorTok{=}\NormalTok{ i }\OperatorTok{*} \DecValTok{4}\OperatorTok{;}
\NormalTok{  out[px]     }\OperatorTok{=}\NormalTok{ c[}\DecValTok{0}\NormalTok{]}\OperatorTok{;}   \CommentTok{// R}
\NormalTok{  out[px }\OperatorTok{+} \DecValTok{1}\NormalTok{] }\OperatorTok{=}\NormalTok{ c[}\DecValTok{1}\NormalTok{]}\OperatorTok{;}   \CommentTok{// G}
\NormalTok{  out[px }\OperatorTok{+} \DecValTok{2}\NormalTok{] }\OperatorTok{=}\NormalTok{ c[}\DecValTok{2}\NormalTok{]}\OperatorTok{;}   \CommentTok{// B}
\NormalTok{  out[px }\OperatorTok{+} \DecValTok{3}\NormalTok{] }\OperatorTok{=} \DecValTok{255}\OperatorTok{;}    \CommentTok{// A (opaque)}
\NormalTok{\}}
\end{Highlighting}
\end{Shaded}

This pass is extremely fast: a single linear scan with array-index
lookups. For a 1920×1080 image, it processes \textasciitilde2 million
pixels in milliseconds.

\subsubsection{Pass 2: Fill UNMATCHED
Pixels}\label{pass-2-fill-unmatched-pixels}

Some pixels remain unpainted after Pass 1. These are pixels that belong
to thin borders, text anti-aliasing fringes, or tiny UI details that
fell below the green-probe threshold. Rather than leaving them
transparent (which would show as black or checkerboard), we fill them by
averaging the colors of their 8-connected neighbors.

\begin{Shaded}
\begin{Highlighting}[]
\CommentTok{// Pass 2: fill UNMATCHED pixels by averaging 8{-}connected assigned neighbors}
\KeywordTok{const}\NormalTok{ fallback }\OperatorTok{=}\NormalTok{ colors[}\StringTok{\textquotesingle{}Main\_Back\textquotesingle{}}\NormalTok{] }\OperatorTok{||}\NormalTok{ [}\DecValTok{20}\OperatorTok{,} \DecValTok{20}\OperatorTok{,} \DecValTok{30}\NormalTok{]}\OperatorTok{;}
\ControlFlowTok{for}\NormalTok{ (}\KeywordTok{let}\NormalTok{ y }\OperatorTok{=} \DecValTok{0}\OperatorTok{;}\NormalTok{ y }\OperatorTok{\textless{}}\NormalTok{ H}\OperatorTok{;}\NormalTok{ y}\OperatorTok{++}\NormalTok{) \{}
  \ControlFlowTok{for}\NormalTok{ (}\KeywordTok{let}\NormalTok{ x }\OperatorTok{=} \DecValTok{0}\OperatorTok{;}\NormalTok{ x }\OperatorTok{\textless{}}\NormalTok{ W}\OperatorTok{;}\NormalTok{ x}\OperatorTok{++}\NormalTok{) \{}
    \KeywordTok{const}\NormalTok{ i }\OperatorTok{=}\NormalTok{ y }\OperatorTok{*}\NormalTok{ W }\OperatorTok{+}\NormalTok{ x}\OperatorTok{;}
    \ControlFlowTok{if}\NormalTok{ (out[i }\OperatorTok{*} \DecValTok{4} \OperatorTok{+} \DecValTok{3}\NormalTok{] }\OperatorTok{\textgreater{}} \DecValTok{0}\NormalTok{) }\ControlFlowTok{continue}\OperatorTok{;} \CommentTok{// already painted}

    \KeywordTok{let}\NormalTok{ rSum }\OperatorTok{=} \DecValTok{0}\OperatorTok{,}\NormalTok{ gSum }\OperatorTok{=} \DecValTok{0}\OperatorTok{,}\NormalTok{ bSum }\OperatorTok{=} \DecValTok{0}\OperatorTok{,}\NormalTok{ count }\OperatorTok{=} \DecValTok{0}\OperatorTok{;}
    \ControlFlowTok{for}\NormalTok{ (}\KeywordTok{let}\NormalTok{ dy }\OperatorTok{=} \OperatorTok{{-}}\DecValTok{1}\OperatorTok{;}\NormalTok{ dy }\OperatorTok{\textless{}=} \DecValTok{1}\OperatorTok{;}\NormalTok{ dy}\OperatorTok{++}\NormalTok{) \{}
      \ControlFlowTok{for}\NormalTok{ (}\KeywordTok{let}\NormalTok{ dx }\OperatorTok{=} \OperatorTok{{-}}\DecValTok{1}\OperatorTok{;}\NormalTok{ dx }\OperatorTok{\textless{}=} \DecValTok{1}\OperatorTok{;}\NormalTok{ dx}\OperatorTok{++}\NormalTok{) \{}
        \ControlFlowTok{if}\NormalTok{ (dy }\OperatorTok{===} \DecValTok{0} \OperatorTok{\&\&}\NormalTok{ dx }\OperatorTok{===} \DecValTok{0}\NormalTok{) }\ControlFlowTok{continue}\OperatorTok{;}
        \KeywordTok{const}\NormalTok{ ny }\OperatorTok{=}\NormalTok{ y }\OperatorTok{+}\NormalTok{ dy}\OperatorTok{,}\NormalTok{ nx }\OperatorTok{=}\NormalTok{ x }\OperatorTok{+}\NormalTok{ dx}\OperatorTok{;}
        \ControlFlowTok{if}\NormalTok{ (ny }\OperatorTok{\textless{}} \DecValTok{0} \OperatorTok{||}\NormalTok{ ny }\OperatorTok{\textgreater{}=}\NormalTok{ H }\OperatorTok{||}\NormalTok{ nx }\OperatorTok{\textless{}} \DecValTok{0} \OperatorTok{||}\NormalTok{ nx }\OperatorTok{\textgreater{}=}\NormalTok{ W) }\ControlFlowTok{continue}\OperatorTok{;}
        \KeywordTok{const}\NormalTok{ ni }\OperatorTok{=}\NormalTok{ ny }\OperatorTok{*}\NormalTok{ W }\OperatorTok{+}\NormalTok{ nx}\OperatorTok{;}
        \KeywordTok{const}\NormalTok{ nIdx }\OperatorTok{=}\NormalTok{ map[ni]}\OperatorTok{;}         \CommentTok{// read from IMMUTABLE map}
        \ControlFlowTok{if}\NormalTok{ (nIdx }\OperatorTok{===}\NormalTok{ UNMATCHED) }\ControlFlowTok{continue}\OperatorTok{;}
        \KeywordTok{const}\NormalTok{ c }\OperatorTok{=}\NormalTok{ colorByIndex[nIdx]}\OperatorTok{;} \CommentTok{// O(1) lookup}
        \ControlFlowTok{if}\NormalTok{ (}\OperatorTok{!}\NormalTok{c) }\ControlFlowTok{continue}\OperatorTok{;}
\NormalTok{        rSum }\OperatorTok{+=}\NormalTok{ c[}\DecValTok{0}\NormalTok{]}\OperatorTok{;}\NormalTok{ gSum }\OperatorTok{+=}\NormalTok{ c[}\DecValTok{1}\NormalTok{]}\OperatorTok{;}\NormalTok{ bSum }\OperatorTok{+=}\NormalTok{ c[}\DecValTok{2}\NormalTok{]}\OperatorTok{;}\NormalTok{ count}\OperatorTok{++;}
\NormalTok{      \}}
\NormalTok{    \}}

    \KeywordTok{const}\NormalTok{ px }\OperatorTok{=}\NormalTok{ i }\OperatorTok{*} \DecValTok{4}\OperatorTok{;}
    \ControlFlowTok{if}\NormalTok{ (count }\OperatorTok{\textgreater{}} \DecValTok{0}\NormalTok{) \{}
\NormalTok{      out[px]     }\OperatorTok{=}\NormalTok{ (rSum }\OperatorTok{/}\NormalTok{ count) }\OperatorTok{|} \DecValTok{0}\OperatorTok{;}
\NormalTok{      out[px }\OperatorTok{+} \DecValTok{1}\NormalTok{] }\OperatorTok{=}\NormalTok{ (gSum }\OperatorTok{/}\NormalTok{ count) }\OperatorTok{|} \DecValTok{0}\OperatorTok{;}
\NormalTok{      out[px }\OperatorTok{+} \DecValTok{2}\NormalTok{] }\OperatorTok{=}\NormalTok{ (bSum }\OperatorTok{/}\NormalTok{ count) }\OperatorTok{|} \DecValTok{0}\OperatorTok{;}
\NormalTok{      out[px }\OperatorTok{+} \DecValTok{3}\NormalTok{] }\OperatorTok{=} \DecValTok{255}\OperatorTok{;}
\NormalTok{    \} }\ControlFlowTok{else}\NormalTok{ \{}
      \CommentTok{// Isolated unmatched region — use background fallback}
\NormalTok{      out[px]     }\OperatorTok{=}\NormalTok{ fallback[}\DecValTok{0}\NormalTok{]}\OperatorTok{;}
\NormalTok{      out[px }\OperatorTok{+} \DecValTok{1}\NormalTok{] }\OperatorTok{=}\NormalTok{ fallback[}\DecValTok{1}\NormalTok{]}\OperatorTok{;}
\NormalTok{      out[px }\OperatorTok{+} \DecValTok{2}\NormalTok{] }\OperatorTok{=}\NormalTok{ fallback[}\DecValTok{2}\NormalTok{]}\OperatorTok{;}
\NormalTok{      out[px }\OperatorTok{+} \DecValTok{3}\NormalTok{] }\OperatorTok{=} \DecValTok{255}\OperatorTok{;}
\NormalTok{    \}}
\NormalTok{  \}}
\NormalTok{\}}
\end{Highlighting}
\end{Shaded}

\textbf{Critical design choice:} We read neighbor indices from the
\textbf{immutable pixel map} (\texttt{map{[}ni{]}}), not from the output
buffer (\texttt{out{[}{]}}). If we read from \texttt{out{[}{]}}, a
filled pixel would influence its own neighbors in the same pass, causing
color to propagate like a flood fill across the entire image. By reading
from the immutable source, each pixel's fill color is determined solely
by the original Pass-1 assignments, producing smooth but bounded
boundary interpolation.

\subsubsection{Pass 3: Composite Text and
Icons}\label{pass-3-composite-text-and-icons}

Pass 1 and Pass 2 produce a clean recolored UI, but text and icons are
still missing. Text pixels were mostly UNMATCHED because the green-probe
method cannot reliably distinguish between ``a green pixel that belongs
to a text character'' and ``a green pixel that belongs to the background
behind the text.'' Anti-aliasing makes this especially hard: the edge of
a letter is a blend of text color and background color.

Our solution is a \textbf{text masking theme}.

\subsection{4.5 Text Masking Theme}\label{text-masking-theme}

We create a special \texttt{.xrnc} theme where \textbf{all UI chrome is
pure black} and \textbf{all text and icons are pure white}. We
screenshot Renoise with this theme loaded. The result is a
black-and-white image where white pixels represent text and icons, and
black pixels represent everything else.

In Pass 3, we use this B/W screenshot as an alpha mask:

\begin{Shaded}
\begin{Highlighting}[]
\CommentTok{// Pass 3: composite text using text{-}theme screenshot as anti{-}aliasing guide}
\ControlFlowTok{if}\NormalTok{ (textThemeData) \{}
  \KeywordTok{const}\NormalTok{ mainFont }\OperatorTok{=}\NormalTok{ colors[}\StringTok{\textquotesingle{}Main\_Font\textquotesingle{}}\NormalTok{] }\OperatorTok{||}\NormalTok{ [}\DecValTok{220}\OperatorTok{,} \DecValTok{220}\OperatorTok{,} \DecValTok{220}\NormalTok{]}\OperatorTok{;}
  \ControlFlowTok{for}\NormalTok{ (}\KeywordTok{let}\NormalTok{ i }\OperatorTok{=} \DecValTok{0}\OperatorTok{;}\NormalTok{ i }\OperatorTok{\textless{}}\NormalTok{ W }\OperatorTok{*}\NormalTok{ H}\OperatorTok{;}\NormalTok{ i}\OperatorTok{++}\NormalTok{) \{}
    \KeywordTok{const}\NormalTok{ px }\OperatorTok{=}\NormalTok{ i }\OperatorTok{*} \DecValTok{4}\OperatorTok{;}
    \CommentTok{// Rec.601 perceptual luminance — human eyes are most sensitive to green}
    \KeywordTok{const}\NormalTok{ lum }\OperatorTok{=} \FloatTok{0.299} \OperatorTok{*}\NormalTok{ textThemeData[px]}
              \OperatorTok{+} \FloatTok{0.587} \OperatorTok{*}\NormalTok{ textThemeData[px }\OperatorTok{+} \DecValTok{1}\NormalTok{]}
              \OperatorTok{+} \FloatTok{0.114} \OperatorTok{*}\NormalTok{ textThemeData[px }\OperatorTok{+} \DecValTok{2}\NormalTok{]}\OperatorTok{;}
    \KeywordTok{const}\NormalTok{ alpha }\OperatorTok{=} \DecValTok{1} \OperatorTok{{-}}\NormalTok{ lum }\OperatorTok{/} \DecValTok{255}\OperatorTok{;} \CommentTok{// black = 0 (UI), white = 1 (text)}
    \ControlFlowTok{if}\NormalTok{ (alpha }\OperatorTok{\textless{}} \FloatTok{0.05}\NormalTok{) }\ControlFlowTok{continue}\OperatorTok{;}  \CommentTok{// skip near{-}black pixels}

    \KeywordTok{const}\NormalTok{ idx }\OperatorTok{=}\NormalTok{ map[i]}\OperatorTok{;}
    \CommentTok{// Skip text compositing on known UI chrome (scrollbars, sliders, etc.)}
    \ControlFlowTok{if}\NormalTok{ (idx }\OperatorTok{!==}\NormalTok{ UNMATCHED }\OperatorTok{\&\&}\NormalTok{ idx }\OperatorTok{\textless{}}\NormalTok{ elements}\OperatorTok{.}\AttributeTok{length} \OperatorTok{\&\&}\NormalTok{ isUIChrome[idx]) }\ControlFlowTok{continue}\OperatorTok{;}

    \KeywordTok{let}\NormalTok{ fontColor }\OperatorTok{=}\NormalTok{ mainFont}\OperatorTok{;}
    \ControlFlowTok{if}\NormalTok{ (idx }\OperatorTok{!==}\NormalTok{ UNMATCHED }\OperatorTok{\&\&}\NormalTok{ idx }\OperatorTok{\textless{}}\NormalTok{ elements}\OperatorTok{.}\AttributeTok{length}\NormalTok{) \{}
      \KeywordTok{const}\NormalTok{ c }\OperatorTok{=}\NormalTok{ colorByIndex[idx]}\OperatorTok{;}
      \ControlFlowTok{if}\NormalTok{ (c) fontColor }\OperatorTok{=}\NormalTok{ c}\OperatorTok{;} \CommentTok{// use element\textquotesingle{}s own font color if available}
\NormalTok{    \}}

    \CommentTok{// Standard alpha blend: fontColor over existing pixel}
\NormalTok{    out[px]     }\OperatorTok{=}\NormalTok{ (fontColor[}\DecValTok{0}\NormalTok{] }\OperatorTok{*}\NormalTok{ alpha }\OperatorTok{+}\NormalTok{ out[px]     }\OperatorTok{*}\NormalTok{ (}\DecValTok{1} \OperatorTok{{-}}\NormalTok{ alpha)) }\OperatorTok{|} \DecValTok{0}\OperatorTok{;}
\NormalTok{    out[px }\OperatorTok{+} \DecValTok{1}\NormalTok{] }\OperatorTok{=}\NormalTok{ (fontColor[}\DecValTok{1}\NormalTok{] }\OperatorTok{*}\NormalTok{ alpha }\OperatorTok{+}\NormalTok{ out[px }\OperatorTok{+} \DecValTok{1}\NormalTok{] }\OperatorTok{*}\NormalTok{ (}\DecValTok{1} \OperatorTok{{-}}\NormalTok{ alpha)) }\OperatorTok{|} \DecValTok{0}\OperatorTok{;}
\NormalTok{    out[px }\OperatorTok{+} \DecValTok{2}\NormalTok{] }\OperatorTok{=}\NormalTok{ (fontColor[}\DecValTok{2}\NormalTok{] }\OperatorTok{*}\NormalTok{ alpha }\OperatorTok{+}\NormalTok{ out[px }\OperatorTok{+} \DecValTok{2}\NormalTok{] }\OperatorTok{*}\NormalTok{ (}\DecValTok{1} \OperatorTok{{-}}\NormalTok{ alpha)) }\OperatorTok{|} \DecValTok{0}\OperatorTok{;}
\NormalTok{    out[px }\OperatorTok{+} \DecValTok{3}\NormalTok{] }\OperatorTok{=} \DecValTok{255}\OperatorTok{;}
\NormalTok{  \}}
\NormalTok{\}}
\end{Highlighting}
\end{Shaded}

\textbf{Why Rec.601 luminance?} Human eyes are roughly 2.9× more
sensitive to green than to blue and 1.7× more sensitive to green than to
red. The Rec.601 formula (\texttt{0.299R\ +\ 0.587G\ +\ 0.114B})
captures this perceptual non-uniformity. A pixel that looks ``medium
gray'' to a human will produce a luminance value near 128, which
translates to an alpha of 0.5 --- exactly what we want for anti-aliased
text edges.

\textbf{Why skip UI chrome elements?} Some UI elements like scrollbars,
sliders, VU meters, and cursors have fine detail that the text-theme
might misidentify as text. We maintain a list of
\texttt{NON\_TEXT\_ELEMENT\_PATTERNS} (\texttt{Scroll}, \texttt{Slider},
\texttt{Border}, \texttt{VU}, \texttt{Meter}, etc.) and skip text
compositing on pixels mapped to those elements. This prevents, for
example, a scrollbar thumb from being recolored as if it were text.

\begin{center}\rule{0.5\linewidth}{0.5pt}\end{center}

\section{5. The Theme Creator Tool}\label{the-theme-creator-tool}

Beyond passive preview, the system includes an active \textbf{Theme
Creator} at \texttt{/create}. This page presents 70 HTML5 color pickers
--- one per Renoise element --- ranked by their pixel coverage.

\subsection{5.1 Coverage Ranking}\label{coverage-ranking}

The color pickers are not listed alphabetically. They are sorted by how
much screen real estate each element controls, computed from the pixel
maps:

\begin{longtable}[]{@{}lll@{}}
\toprule\noalign{}
Rank & Element & Coverage \\
\midrule\noalign{}
\endhead
\bottomrule\noalign{}
\endlastfoot
1 & \texttt{Main\_Back} & 46.38\% \\
2 & \texttt{Body\_Back} & 31.50\% \\
3 & \texttt{Button\_Back} & 6.61\% \\
4 & \texttt{ValueBox\_Back} & 4.10\% \\
5 & \texttt{Pattern\_Default\_Back} & 3.51\% \\
6 & \texttt{Main\_Font} & 2.53\% \\
\ldots{} & \ldots{} & \ldots{} \\
70 & \texttt{VuMeter\_Meter\_High} & \textasciitilde0.00\% \\
\end{longtable}

This ranking is load-bearing for usability. If \texttt{Main\_Back} and
\texttt{Body\_Back} were buried at the bottom of the list, users would
need to scroll through 50+ obscure elements before changing the two
colors that define 78\% of the theme's appearance. By surfacing
high-coverage elements first, a user can create a recognizable theme in
under 30 seconds.

\subsection{5.2 Live Preview and
Download}\label{live-preview-and-download}

The creator page communicates with two API endpoints:

\begin{itemize}
\tightlist
\item
  \textbf{\texttt{POST\ /api/render-preview}} --- Accepts a JSON
  \texttt{elementColorMap}, normalizes hex strings to
  \texttt{{[}r,g,b{]}} arrays, calls \texttt{generatePreviews()}, and
  returns preview image URLs.
\item
  \textbf{\texttt{POST\ /api/download-xrnc}} --- Accepts the same color
  map, calls \texttt{generateXrnc()} to produce a valid \texttt{.xrnc}
  XML string, and streams it as a file download.
\end{itemize}

Both endpoints are protected by CSRF tokens (see Section 6). The preview
endpoint returns a \texttt{previewSlug} that identifies the temporary
preview directory; images are served from
\texttt{public/uploads/previews/} and cleaned up by standard server
maintenance.

\begin{center}\rule{0.5\linewidth}{0.5pt}\end{center}

\section{6. Security \& Hardening}\label{security-hardening}

A theme gallery is a public-facing file upload service. We implement
defense in depth across multiple layers.

\subsection{6.1 CSRF Protection}\label{csrf-protection}

All mutating routes (upload, comment, login, register, like, preview
render, XRNC download) use an \textbf{HMAC-based double-submit cookie
pattern}:

\begin{enumerate}
\def\labelenumi{\arabic{enumi}.}
\tightlist
\item
  Each session receives a random 256-bit \texttt{csrfSecret} stored
  server-side.
\item
  Per-request tokens are computed as
  \texttt{HMAC-SHA256(csrfSecret,\ "csrf:"\ +\ hourlyTimestamp)}.
\item
  Tokens rotate hourly. The verification window accepts the current or
  previous hour's token to handle clock skew and in-flight requests.
\item
  Comparison uses \texttt{crypto.timingSafeEqual()} to prevent timing
  attacks.
\item
  AJAX API requests are exempted via the
  \texttt{X-Requested-With:\ XMLHttpRequest} header, which browsers
  enforce as a same-origin constraint.
\end{enumerate}

\subsection{6.2 Session Hardening}\label{session-hardening}

Sessions are managed by \texttt{express-session} with the following
flags: - \texttt{httpOnly:\ true} --- Prevents JavaScript from reading
the session cookie -
\texttt{sameSite:\ \textquotesingle{}lax\textquotesingle{}} --- Prevents
cross-origin POST attacks while allowing normal navigation -
\texttt{secure:\ true} --- Enforces HTTPS in production -
\texttt{maxAge:\ 30\ days} --- Balances convenience with security

\subsection{6.3 Rate Limiting}\label{rate-limiting}

\begin{itemize}
\tightlist
\item
  \textbf{Auth routes} (\texttt{/login}, \texttt{/register},
  \texttt{/forgot-password}, \texttt{/reset-password}): 12 requests per
  15-minute window per IP.
\item
  \textbf{Download endpoint} (\texttt{/download/:slug}): 30 requests per
  minute per IP.
\end{itemize}

\subsection{6.4 Input Sanitization}\label{input-sanitization}

\begin{itemize}
\tightlist
\item
  \textbf{Path traversal:} All user-supplied filenames are sanitized
  with \texttt{path.basename()} before storage.
\item
  \textbf{File types:} Multer's \texttt{fileFilter} rejects
  non-\texttt{.xrnc}/\texttt{.xml} theme uploads and non-image
  screenshot uploads.
\item
  \textbf{File size:} Uploads are capped at 10 MB.
\item
  \textbf{XML parsing:} The XML parser disables entity expansion
  (\texttt{processEntities:\ false}) and ignores declarations/processing
  instructions, preventing XXE (XML External Entity) attacks.
\end{itemize}

\subsection{6.5 Password Security}\label{password-security}

User passwords are hashed with \textbf{bcrypt at 12 rounds}. The
registration form enforces a minimum length of 8 characters. Password
reset tokens are cryptographically random 256-bit values with a 1-hour
expiration.

\begin{center}\rule{0.5\linewidth}{0.5pt}\end{center}

\section{7. Headless Automation
Pipeline}\label{headless-automation-pipeline}

Generating 71 screenshots manually would take hours and introduce human
error (window positioning, timing, song state). We automated the entire
process using a Linux headless stack.

\subsection{7.1 Tools Used}\label{tools-used}

\begin{longtable}[]{@{}
  >{\raggedright\arraybackslash}p{(\linewidth - 2\tabcolsep) * \real{0.5000}}
  >{\raggedright\arraybackslash}p{(\linewidth - 2\tabcolsep) * \real{0.5000}}@{}}
\toprule\noalign{}
\begin{minipage}[b]{\linewidth}\raggedright
Tool
\end{minipage} & \begin{minipage}[b]{\linewidth}\raggedright
Role
\end{minipage} \\
\midrule\noalign{}
\endhead
\bottomrule\noalign{}
\endlastfoot
\textbf{Xvfb} & Virtual X11 display (1920×1080×24) --- no physical
monitor needed \\
\textbf{xdotool} & Window detection and geometry verification \\
\textbf{maim} & Screenshot capture by window ID \\
\textbf{Renoise} & The DAW itself, launched with overridden config \\
\textbf{Python 3} & Config injection via regex substitution \\
\textbf{Bash} & Orchestration loop with cleanup traps \\
\end{longtable}

\subsection{7.2 Temp Config Isolation}\label{temp-config-isolation}

The script \textbf{never touches the user's real Renoise config}. It
creates a temporary directory structure and overrides \texttt{HOME} so
Renoise reads from \texttt{/tmp/rns-capture-cfg/} instead of
\texttt{\textasciitilde{}/.config/Renoise/}.

\begin{Shaded}
\begin{Highlighting}[]
\VariableTok{TEMP\_CFG}\OperatorTok{=}\StringTok{"/tmp/rns{-}capture{-}cfg/V3.5.4"}
\VariableTok{HOME\_OVERRIDE}\OperatorTok{=}\StringTok{"/tmp/rns{-}capture{-}home"}
\FunctionTok{mkdir} \AttributeTok{{-}p} \StringTok{"}\VariableTok{$TEMP\_CFG}\StringTok{"}
\FunctionTok{cp} \StringTok{"}\VariableTok{$REAL\_CONFIG}\StringTok{"} \StringTok{"}\VariableTok{$TEMP\_CFG}\StringTok{/Config.xml"}
\CommentTok{\# Plugin caches copied for fast startup}
\FunctionTok{cp} \StringTok{"}\VariableTok{$HOME}\StringTok{/.config/Renoise/V3.5.4/"}\PreprocessorTok{*}\NormalTok{.db }\StringTok{"}\VariableTok{$TEMP\_CFG}\StringTok{/"}
\FunctionTok{ln} \AttributeTok{{-}sf} \StringTok{"}\VariableTok{$TEMP\_CFG}\StringTok{"} \StringTok{"}\VariableTok{$HOME\_OVERRIDE}\StringTok{/.config/Renoise/V3.5.4"}
\end{Highlighting}
\end{Shaded}

\subsection{7.3 Renoise Launch
Optimization}\label{renoise-launch-optimization}

Renoise normally rescans plugins on startup, which can take 30+ seconds.
The script patches the config XML to disable all rescanning:

\begin{Shaded}
\begin{Highlighting}[]
\NormalTok{c }\OperatorTok{=}\NormalTok{ re.sub(}\VerbatimStringTok{r\textquotesingle{}\textless{}RescanPreviouslyFailedPlugs\textgreater{}.*\textless{}/RescanPreviouslyFailedPlugs\textgreater{}\textquotesingle{}}\NormalTok{,}
           \StringTok{\textquotesingle{}\textless{}RescanPreviouslyFailedPlugs\textgreater{}false\textless{}/RescanPreviouslyFailedPlugs\textgreater{}\textquotesingle{}}\NormalTok{, c)}
\NormalTok{c }\OperatorTok{=}\NormalTok{ re.sub(}\VerbatimStringTok{r\textquotesingle{}\textless{}ScanForNewPluginsOnStartup\textgreater{}.*\textless{}/ScanForNewPluginsOnStartup\textgreater{}\textquotesingle{}}\NormalTok{,}
           \StringTok{\textquotesingle{}\textless{}ScanForNewPluginsOnStartup\textgreater{}false\textless{}/ScanForNewPluginsOnStartup\textgreater{}\textquotesingle{}}\NormalTok{, c)}
\NormalTok{c }\OperatorTok{=}\NormalTok{ re.sub(}\VerbatimStringTok{r\textquotesingle{}\textless{}AutoRescanHotPluggedDevices\textgreater{}.*\textless{}/AutoRescanHotPluggedDevices\textgreater{}\textquotesingle{}}\NormalTok{,}
           \StringTok{\textquotesingle{}\textless{}AutoRescanHotPluggedDevices\textgreater{}false\textless{}/AutoRescanHotPluggedDevices\textgreater{}\textquotesingle{}}\NormalTok{, c)}
\end{Highlighting}
\end{Shaded}

It also forces windowed mode to 1920×1080 and maximizes the window,
ensuring consistent screenshot dimensions.

\subsection{7.4 Fast Kill-Restart Cycle}\label{fast-kill-restart-cycle}

Instead of gracefully shutting down Renoise (which triggers slow JACK
teardown and config writes), the script uses \texttt{kill\ -9} after
each screenshot:

\begin{Shaded}
\begin{Highlighting}[]
\BuiltInTok{kill} \AttributeTok{{-}9} \VariableTok{$RPID} \DecValTok{2}\OperatorTok{\textgreater{}}\NormalTok{/dev/null }\KeywordTok{||} \FunctionTok{true}
\BuiltInTok{wait} \VariableTok{$RPID} \DecValTok{2}\OperatorTok{\textgreater{}}\NormalTok{/dev/null }\KeywordTok{||} \FunctionTok{true}
\FunctionTok{sleep}\NormalTok{ 0.5}
\end{Highlighting}
\end{Shaded}

This reduces the per-variant cycle to approximately \textbf{1 second of
Renoise runtime}. The full 71-variant capture completes in roughly
\textbf{3 minutes}. Before each launch, the script restores the clean
baseline config so that variants do not leak state into one another.

\subsection{7.5 Capture Verification}\label{capture-verification}

The script validates each screenshot by file size. Screenshots under 10
KB are rejected as failures (usually indicating a window detection
timeout or crash). Failed variants are logged but do not abort the
pipeline.

\begin{center}\rule{0.5\linewidth}{0.5pt}\end{center}

\section{8. Results \& Impact}\label{results-impact}

\subsection{8.1 Map Coverage}\label{map-coverage}

The current pixel maps cover \textbf{74 Renoise UI elements} across
three views:

\begin{longtable}[]{@{}lll@{}}
\toprule\noalign{}
View & Resolution & Elements Mapped \\
\midrule\noalign{}
\endhead
\bottomrule\noalign{}
\endlastfoot
Pattern Editor & 1920×1080 & 74 \\
Mixer & 1920×1080 & 74 (shared map) \\
Waveform & 1920×1080 & 74 (shared map) \\
\end{longtable}

Total matched pixels across all three views: \textbf{4,868,286} (from
the coverage analysis used to rank the Theme Creator).

The top coverage elements dominate the visual identity of any theme:

\begin{verbatim}
01. Main_Back                    46.38%  (892,348 px)
02. Body_Back                    31.50%  (606,042 px)
03. Button_Back                   6.61%  (127,210 px)
04. ValueBox_Back                 4.10%   (78,902 px)
05. Pattern_Default_Back          3.51%   (67,574 px)
06. Main_Font                     2.53%   (48,728 px)
07. Scrollbar                     1.32%   (25,416 px)
08. Selected_Button_Back          0.99%   (19,068 px)
09. StandBy_Selection_Back        0.66%   (12,712 px)
10. Pattern_CenterBar_Back        0.40%    (7,704 px)
\end{verbatim}

These 10 elements control \textbf{97.9\%} of the screen. A theme author
who only adjusts these 10 colors can create a preview that looks
substantially different from the default --- proof that the coverage
ranking in the Theme Creator is not merely informative but functionally
critical.

\subsection{8.2 Preview Accuracy}\label{preview-accuracy}

The three-pass renderer produces images that are \textbf{pixel-perfect
in color assignment} for all mapped elements. UNMATCHED pixels (text
fringes, thin borders, anti-aliasing) are filled via neighbor averaging
in Pass 2 and then corrected via text-theme compositing in Pass 3. The
result is visually indistinguishable from a native Renoise screenshot at
normal viewing distance.

The system achieves this accuracy \textbf{without any Renoise code,
assets, or runtime dependency} on the server. The only Renoise
dependency is during the one-time map generation phase.

\begin{center}\rule{0.5\linewidth}{0.5pt}\end{center}

\section{9. Future Work}\label{future-work}

\subsection{9.1 Per-View Map Refinement}\label{per-view-map-refinement}

Currently, Pattern Editor, Mixer, and Waveform views share the same
element list but use separate binary map files (\texttt{pattern.bin},
\texttt{mixer.bin}, \texttt{waveform.bin}). Future work includes
refining each view's map independently, since some elements (e.g.,
\texttt{VuMeter\_Meter}) appear in different positions and proportions
across views.

\subsection{9.2 SVG-Based Vector Preview
Renderer}\label{svg-based-vector-preview-renderer}

The map generation pipeline already outputs SVG vector outlines of
element regions (\texttt{maps/pattern.svg}, \texttt{maps/mixer.svg},
\texttt{maps/waveform.svg}). A future renderer could use these SVGs
instead of per-pixel binary maps, enabling: - \textbf{Arbitrary
resolution scaling} ( previews at 4K or mobile sizes without pixelation)
- \textbf{Smaller file sizes} (SVG paths compress better than raw pixel
arrays for large uniform regions) - \textbf{Interactive hover effects}
(highlighting which UI element corresponds to which color definition)

\subsection{9.3 Multi-View Capture per
Launch}\label{multi-view-capture-per-launch}

The current automation captures one view per Renoise launch. A single
launch could theoretically capture all three views by programmatically
switching tabs (via xdotool key events) before taking three screenshots.
This would reduce the 71-variant × 3-view capture from
\textasciitilde213 launches to \textasciitilde71 launches --- a 3×
speedup in map maintenance.

\subsection{9.4 Theme Diff and Merge
Tools}\label{theme-diff-and-merge-tools}

With structured \texttt{elementColorMap} objects, we can compute the
perceptual distance between two themes. Future features could include: -
``Find themes similar to this one'' - ``Show me what changed between
version 1 and 2'' - ``Blend two themes'' (interpolate each element color
independently)

\begin{center}\rule{0.5\linewidth}{0.5pt}\end{center}

\section{10. References \& File Index}\label{references-file-index}

\subsection{Core Libraries}\label{core-libraries}

\begin{longtable}[]{@{}
  >{\raggedright\arraybackslash}p{(\linewidth - 2\tabcolsep) * \real{0.4000}}
  >{\raggedright\arraybackslash}p{(\linewidth - 2\tabcolsep) * \real{0.6000}}@{}}
\toprule\noalign{}
\begin{minipage}[b]{\linewidth}\raggedright
File
\end{minipage} & \begin{minipage}[b]{\linewidth}\raggedright
Purpose
\end{minipage} \\
\midrule\noalign{}
\endhead
\bottomrule\noalign{}
\endlastfoot
\texttt{lib/parser.js} & Parses \texttt{.xrnc} XML; extracts colors;
assigns semantic roles and weights; exports \texttt{elementColorMap} \\
\texttt{lib/categorize.js} & Converts colors to HSL; applies weighted
statistical analysis to auto-tag themes (dark/light, warm/cool, etc.) \\
\texttt{lib/palette.js} & Generates weighted SVG palette strips with
MAIN/SECONDARY/UI/ACCENTS tiers \\
\texttt{lib/preview-renderer.js} & Three-pass pixel map renderer: paint
→ fill → text composite \\
\texttt{lib/database.js} & SQLite schema, prepared statements, batched
metadata attachment, authentication \\
\texttt{lib/xrnc-generator.js} & Generates valid \texttt{.xrnc} XML from
an \texttt{elementColorMap} for the Theme Creator download feature \\
\end{longtable}

\subsection{Automation Tools}\label{automation-tools}

\begin{longtable}[]{@{}
  >{\raggedright\arraybackslash}p{(\linewidth - 2\tabcolsep) * \real{0.4000}}
  >{\raggedright\arraybackslash}p{(\linewidth - 2\tabcolsep) * \real{0.6000}}@{}}
\toprule\noalign{}
\begin{minipage}[b]{\linewidth}\raggedright
File
\end{minipage} & \begin{minipage}[b]{\linewidth}\raggedright
Purpose
\end{minipage} \\
\midrule\noalign{}
\endhead
\bottomrule\noalign{}
\endlastfoot
\texttt{tools/generate-white-variants.js} & Generates 71 variant themes:
1 white baseline + 70 green-probe variants \\
\texttt{tools/capture-variants.sh} & Headless bash pipeline: Xvfb →
Renoise → maim screenshot → kill-9 cycle \\
\texttt{tools/build-diff-maps.js} & Builds \texttt{maps/pattern.bin} and
\texttt{maps/pattern.json} from probe screenshots via green-delta
detection and overlap resolution \\
\texttt{tools/generate-diff-variants.js} & Alternative variant generator
using magenta probes against a non-white baseline
(legacy/supplemental) \\
\end{longtable}

\subsection{Map Assets}\label{map-assets}

\begin{longtable}[]{@{}
  >{\raggedright\arraybackslash}p{(\linewidth - 2\tabcolsep) * \real{0.4000}}
  >{\raggedright\arraybackslash}p{(\linewidth - 2\tabcolsep) * \real{0.6000}}@{}}
\toprule\noalign{}
\begin{minipage}[b]{\linewidth}\raggedright
File
\end{minipage} & \begin{minipage}[b]{\linewidth}\raggedright
Purpose
\end{minipage} \\
\midrule\noalign{}
\endhead
\bottomrule\noalign{}
\endlastfoot
\texttt{maps/pattern.bin} & Flat \texttt{Uint8Array} (1920×1080 =
2,073,600 bytes). Each byte = element index or 255 (UNMATCHED) \\
\texttt{maps/pattern.json} & Metadata: width, height, ordered element
name array \\
\texttt{maps/mixer.bin} / \texttt{maps/mixer.json} & Mixer view pixel
maps \\
\texttt{maps/waveform.bin} / \texttt{maps/waveform.json} & Waveform view
pixel maps \\
\texttt{maps/pattern-text.png} & Cached text-theme screenshot for Pass 3
compositing \\
\texttt{maps/*.svg} & Vector outlines of element regions (future SVG
renderer) \\
\end{longtable}

\subsection{Application Entry Point}\label{application-entry-point}

\begin{longtable}[]{@{}
  >{\raggedright\arraybackslash}p{(\linewidth - 2\tabcolsep) * \real{0.4000}}
  >{\raggedright\arraybackslash}p{(\linewidth - 2\tabcolsep) * \real{0.6000}}@{}}
\toprule\noalign{}
\begin{minipage}[b]{\linewidth}\raggedright
File
\end{minipage} & \begin{minipage}[b]{\linewidth}\raggedright
Purpose
\end{minipage} \\
\midrule\noalign{}
\endhead
\bottomrule\noalign{}
\endlastfoot
\texttt{app.js} & Express server: routes, middleware, upload handling,
CSRF protection, rate limiting, session config, Theme Creator API \\
\end{longtable}

\begin{center}\rule{0.5\linewidth}{0.5pt}\end{center}

\section{Acknowledgments}\label{acknowledgments}

The ``Soon Soon'' demo song by \textbf{Hunz} is used during automated
capture to ensure all UI elements (pattern data, VU meters, automation
envelopes, device chains) are visible and mappable. The Renoise team
built an extraordinarily customizable UI; this system exists to make
that customization more accessible to the community.

\begin{center}\rule{0.5\linewidth}{0.5pt}\end{center}

\emph{Document version: 1.0}\\
\emph{Generated for the Renoise Themes open-source project}\\
\emph{For questions or contributions, refer to the file index in Section
10.}

\end{document}
